\documentclass[12pt,a4paper]{article}

% --- Кодировки и язык ---
\usepackage[T2A]{fontenc}
\usepackage[utf8]{inputenc}
\usepackage[english,russian]{babel}
\usepackage{cmap}

% --- Математика ---
\usepackage{amsmath,amssymb,amsthm,mathtools}
\usepackage{bm}

% --- Геометрия и типографика ---
\usepackage[left=2.5cm,right=2cm,top=2.5cm,bottom=2.5cm]{geometry}
\usepackage{setspace}
\onehalfspacing
\usepackage{indentfirst}
\usepackage[protrusion=true,expansion=false]{microtype}

% --- Библиография и ссылки ---
\usepackage[numbers,sort&compress]{natbib}
\usepackage[hidelinks]{hyperref}
\usepackage{url}

% --- Теоремные окружения ---
\newtheorem{theorem}{Теорема}
\newtheorem{proposition}{Предложение}
\newtheorem{lemma}{Лемма}
\newtheorem{corollary}{Следствие}
\theoremstyle{definition}
\newtheorem{definition}{Определение}
\newtheorem{assumption}{Предположение}
\newtheorem{remark}{Замечание}

% --- Удобные команды ---
\newcommand{\R}{\mathbb{R}}
\newcommand{\N}{\mathbb{N}}
\newcommand{\E}{\mathbb{E}}
\renewcommand{\Pr}{\mathbb{P}}
\DeclareMathOperator*{\argmax}{arg\,max}
\DeclareMathOperator*{\argmin}{arg\,min}
\newcommand{\ind}[1]{\mathbf{1}\!\left\{#1\right\}}

\renewcommand{\proofname}{Доказательство}

% --- Метаданные ---
\title{\textbf{Многоагентная теоретико-игровая модель
    \\ стратегического научного публикования \\ с целевой функцией
  h-индекса РИНЦ}\\[0.3em]
  \large A Multi-Agent Game-Theoretic Model of Strategic Scientific
Publishing with the RINTs h-Index as the Objective Function}

\author{М.\,Г.~Смирнов\thanks{e-mail:
  \texttt{hello@maksim-smirnov.dev}}\\[0.3em]
\small Московский физико-технический институт}

\date{Апрель 2026}

\begin{document}

\maketitle

\begin{center}
  \small УДК 519.834 + 001.891.3\\
  {\footnotesize \textit{MSC2020:} 91A80, 91A40, 91A10, 91A26}
\end{center}

% ============================================================
\begin{abstract}
  \noindent
  В работе строится многоагентная байесовская теоретико-игровая
  модель научного публикования, в которой игроками являются
  автор-исследователь, редакторы журналов, рецензенты и потенциальные
  соавторы. Целевая функция автора задана как дисконтированная
  линейная комбинация четырёх вариантов h-индекса Российского индекса
  научного цитирования (РИНЦ): полного, без самоцитирований автора,
  без самоцитирований соавторов и по ядру РИНЦ. Модель объединяет
  пять формальных блоков: сигнальную подыгру выбора журнала, подыгру
  редактор-рецензент с требованием принудительного цитирования
  (coercive citation), кооперативную подмодель соавторства со
  значением Шепли, повторяющуюся игру цитатных картелей
  (\emph{citation rings}) и эволюционную динамику стратегий в
  популяции учёных. Доказывается существование разделяющего
  совершенного байесовского равновесия в подыгре выбора журнала,
  получается верхняя оценка прироста h-индекса от стратегического
  самоцитирования, формулируется критерий устойчивости кооперативной
  коалиции соавторов через непустоту ядра кооперативной игры, а также
  условия существования эволюционно устойчивой стратегии в популяции
  при институциональных параметрах российской научной системы. Модель
  даёт ряд численно проверяемых предсказаний: оптимальная доля
  самоцитирований находится в интервале 15--25\,\%, устойчивый размер
  цитатного картеля ограничен 2--4 участниками, стандартные
  РИНЦ-метрики с равными долями соавторства систематически смещены
  относительно значения Шепли для критических соавторов.

  \smallskip
  \noindent\textit{Ключевые слова:} теория игр, байесовские
  равновесия, значение Шепли, эволюционно устойчивая стратегия,
  наукометрия, h-индекс, РИНЦ, самоцитирование, соавторство, coercive citation.
\end{abstract}

\begin{otherlanguage}{english}
  \begin{abstract}
    \noindent
    We construct a multi-agent Bayesian game-theoretic model of
    scientific publishing with four types of players: the
    author-researcher, journal editors, reviewers, and potential
    co-authors. The author's objective is a discounted linear
    combination of four variants of the RINTs (Russian Index of
    Scientific Citation) h-index: full, author-self-citation-free,
    coauthor-self-citation-free, and the core-RINTs h-index. The
    model integrates five formal blocks: a signaling subgame of
    journal choice, an editor-reviewer subgame with coercive
    citation, a cooperative sub-model of co-authorship via the
    Shapley value, a repeated game of citation rings, and
    evolutionary dynamics of strategies in a population of scholars.
    We prove the existence of a separating perfect Bayesian
    equilibrium in the journal-choice subgame, derive an upper bound
    on the h-index gain from strategic self-citation, formulate a
    stability criterion for co-authorship coalitions via
    non-emptiness of the core, and provide conditions for the
    existence of an evolutionarily stable strategy under the
    institutional parameters of the Russian academic system. The
    model yields several numerically testable predictions: the
    optimal self-citation rate lies in the 15--25\% interval, the
    stable size of a citation ring is bounded by 2--4 participants,
    and standard RINTs-based equal-share co-authorship counters are
    systematically biased relative to the Shapley value for critical co-authors.

    \smallskip
    \noindent\textit{Keywords:} game theory, Bayesian equilibria,
    Shapley value, evolutionarily stable strategy, scientometrics,
    h-index, RINTs, self-citation, co-authorship, coercive citation.
  \end{abstract}
\end{otherlanguage}

\tableofcontents

% ============================================================
\section{Введение}
\label{sec:intro}

Рост значимости формальных наукометрических показателей в системе
оценки научного труда превратил стратегическое поведение
исследователей в самостоятельный объект экономического и
математического анализа. В российской институциональной среде
ключевую роль играют показатели, рассчитываемые по Российскому
индексу научного цитирования: h-индекс автора в нескольких
модификациях, квартильная принадлежность журналов в системе Russian
Science Citation Index (RSCI), принадлежность ядру РИНЦ, категория
журнала в Перечне ВАК (К1, К2, К3) и уровень в Белом списке журналов
РЦНИ~\citep{akoev2021, pislyakov2022}. Эти показатели используются
при защите диссертаций, аттестации научных сотрудников, распределении
грантов РНФ, формировании комплексного балла публикационной
результативности (КБПР) и премиальных выплат в рамках эффективного контракта.

Такая институциональная архитектура порождает сложную многоагентную
систему, в которой действия каждого участника (автора, редактора,
рецензента, соавтора) стратегически взаимозависимы. Автор выбирает
журнал для подачи, структуру списка литературы, число и состав
соавторов, долю самоцитирований. Редактор определяет порог принятия,
пул рецензентов и (в эмпирически документируемых случаях) требует
дополнительных ссылок на свой журнал. Рецензент выбирает уровень
усилий и рекомендацию. Соавтор соглашается или отказывается от
участия в совместной работе. Отдельные аспекты этой системы изучались
в международной литературе~\citep{garcia2015authored,
  heintzelman2009, bartneck2011, smaldino2016, wilhite2012, foster2015,
ioannidis2019}. Однако единой модели, которая одновременно охватывала
бы многоагентную байесовскую структуру, кооперативный аспект
соавторства, повторяющуюся структуру циклов цитирования и специфику
институтов РИНЦ, в литературе нет.

Настоящая работа заполняет эту пробел. Основные результаты таковы:
\begin{itemize}
  \item формулируется композитная задача максимизации
    дисконтированного h-индекса РИНЦ в четырёх вариантах
    (Раздел~\ref{sec:model});
  \item в подыгре выбора журнала доказывается существование
    разделяющего совершенного байесовского равновесия при условии
    single-crossing (Теорема~\ref{thm:pbe-journal});
  \item в подыгре редактор-рецензент описывается равновесие с
    coercive citation и выводится неравенство, определяющее порог
    редакторского требования (Предложение~\ref{prop:coercive});
  \item в кооперативной подмодели формулируется критерий устойчивости
    коалиции соавторов через непустоту ядра игры; уточняется, что
    отклонение распределения кредита от равных долей возникает лишь
    при явном моделировании неравных вкладов внутри публикации, тогда
    как в стандартной постановке равные доли тождественны значению
    Шепли~\citep{karpov2014} (Теорема~\ref{thm:shapley});
  \item для повторяющейся игры цитатных картелей устанавливается
    верхняя граница размера устойчивого картеля, зависящая от
    дисконт-фактора и вероятности разоблачения (Теорема~\ref{thm:folk});
  \item для эволюционной подмодели выводятся условия существования
    ESS при параметрах российской системы (Предложение~\ref{prop:ess});
  \item для оптимального самоцитирования формулируется и доказывается
    верхняя оценка прироста h-индекса, следующая из пороговой
    структуры индекса (Теорема~\ref{thm:selfcite}).
\end{itemize}

Работа построена следующим образом. В Разделе~\ref{sec:model}
вводится общая структура игры, информационное разбиение и целевая
функция автора. Разделы~\ref{sec:journal-choice}--\ref{sec:selfcite}
последовательно анализируют пять формальных блоков модели. В
Разделе~\ref{sec:integration} обсуждается интеграция равновесий в
единую стационарную конфигурацию и её эмпирические импликации.
Раздел~\ref{sec:conclusion} подводит итоги. Эмпирическая проверка
предсказаний модели на данных РИНЦ, Перечня ВАК и Белого списка
выделена во вторую часть исследования.

% ============================================================
\section{Многоагентная байесовская игра: формальная структура}
\label{sec:model}

\subsection{Игроки, типы и информация}

Рассматривается многоагентная байесовская игра
\begin{equation}
  \Gamma^{b} = \bigl\langle N, (A_i)_{i\in N}, (T_i)_{i\in N}, \mu,
  (u_i)_{i\in N} \bigr\rangle,
  \label{eq:game-structure}
\end{equation}
где $N = \{A\} \cup \mathcal{E} \cup \mathcal{R} \cup \mathcal{C}$:
$A$~--- автор-протагонист, $\mathcal{E} = \{E_1, \dots, E_J\}$~---
множество редакторов журналов, $\mathcal{R}$~--- пул рецензентов,
$\mathcal{C}$~--- множество потенциальных соавторов.

Тип автора $\theta_A \in \Theta_A = [0,1]$ трактуется как латентное
качество исследовательской программы и распределён по функции $F$ с
плотностью $f$, имеющей убывающее отношение правдоподобия. Тип
редактора журнала $j$~--- пара $(\kappa_j, \gamma_j) \in [0,1]^2$,
где $\kappa_j$~--- порог принятия (прокси для престижа),
$\gamma_j$~--- склонность к требованию принудительного цитирования.
Тип рецензента $\theta_R \in [0,1]$ характеризует точность его
сигнала о качестве рукописи.

Априорное совместное распределение типов $\mu \in \Delta(\Theta_A
\times \Theta_E^J \times \Theta_R)$ является общим знанием. Типы
игроков реализуются независимо; каждый игрок знает свой тип, но не
знает типов других игроков.

\subsection{Целевая функция автора}

Пусть в момент времени $t \in \{1, 2, \dots, T\}$ автор имеет вектор
цитирований своих публикаций $\bm{c}_t = (c_{1,t}, \dots,
c_{n_t,t})$, упорядоченный по невозрастанию. Классический h-индекс
определяется как
\begin{equation}
  h(\bm{c}) = \max\bigl\{k \in \N : c_k \geq k\bigr\} = \max\bigl\{q
  \in \N : |\{i : c_i \geq q\}| \geq q\bigr\}.
  \label{eq:hindex-def}
\end{equation}
В системе РИНЦ рассчитываются по меньшей мере четыре варианта
h-индекса: полный $h^{\mathrm{full}}$, без самоцитирований автора
$h^{\mathrm{no\text{-}self}}$, без самоцитирований соавторов
$h^{\mathrm{no\text{-}coauth}}$ и по ядру РИНЦ $h^{\mathrm{core}}$.
Различные институциональные процедуры в России используют разные
комбинации этих показателей~\citep{akoev2021}.

Целевая функция автора имеет вид
\begin{equation}
  U_A(\sigma_A; \sigma_{-A}) = \sum_{t=1}^{T} \delta^{t-1} \Bigl[
    \alpha_1 h_t^{\mathrm{full}} + \alpha_2
    h_t^{\mathrm{no\text{-}self}} + \alpha_3
    h_t^{\mathrm{no\text{-}coauth}} + \alpha_4 h_t^{\mathrm{core}}
  \Bigr] - C(\sigma_A) - \lambda \,\E[\mathrm{Pen}(\sigma_A)],
  \label{eq:utility-author}
\end{equation}
где $\delta \in (0,1)$~--- дисконт-фактор, $\alpha_k \geq 0$ с
$\sum_k \alpha_k = 1$~--- весовые коэффициенты, $C(\sigma_A)$~---
совокупные издержки усилий, $\mathrm{Pen}(\sigma_A)$~--- штрафной
член (риск ретракции, исключения журнала, репутационных потерь),
$\lambda > 0$~--- коэффициент отвращения к риску санкций.

\subsection{Стратегическое пространство автора}

Композитная стратегия автора в каждый период
\begin{equation}
  \sigma_A(\theta_A) = \bigl(j_t,\, e_t,\, \bm{r}_t,\, \bm{s}_t,\,
  N_{P,t},\, \tau_t,\, \ell_t\bigr),
  \label{eq:author-strategy}
\end{equation}
где $j_t \in \{1, \dots, J\}$~--- выбор журнала; $e_t \geq 0$~---
уровень исследовательских усилий; $\bm{r}_t \in \N^J$~--- вектор
ссылок на журналы в списке литературы, $R_t = \|\bm{r}_t\|_1$~---
общая длина списка; $\bm{s}_t \in \N^{n_{t-1}}$~--- вектор
самоцитирований по собственным публикациям; $N_{P,t} \subseteq
\mathcal{C} \cup \{A\}$~--- состав соавторской коалиции для
публикации $P_t$; $\tau_t \in \{0,1\}$~--- выбор между традиционной и
инновационной тематикой; $\ell_t \in \{\text{ru}, \text{en},
\text{both}\}$~--- язык публикации.

Стратегия редактора журнала $j$: функция $\sigma_{E_j}\colon \Theta_A
\times \mathcal{M} \to \{\text{desk-reject}, \text{review}\} \times
\N \times \{\text{accept}, \text{reject}\}$, где $\mathcal{M}$~---
множество рукописей; второй компонент кодирует число дополнительных
ссылок на журнал $j$, требуемых редактором.

Стратегия рецензента: пара $(e_R, s_R) \in \R_+ \times \{A, W, R\}$,
где $e_R$~--- усилие рецензирования, $s_R$~--- рекомендация (accept,
weak, reject).

Стратегия соавтора: присоединение к коалиции и уровень вклада в исследование.

% ============================================================
\section{Подыгра выбора журнала: сигнальная игра}
\label{sec:journal-choice}

\subsection{Постановка}

Подыгра выбора журнала $\Gamma^{\mathrm{J}}$~--- сигнальная игра в
духе~\citep{spence1973}. Автор типа $\theta_A$ выбирает журнал $j \in
\{1, \dots, J\}$; журналы упорядочены по невозрастанию престижа $q_1
> q_2 > \dots > q_J$ с соответствующими порогами принятия $\kappa_1 >
\kappa_2 > \dots > \kappa_J$. Редактор журнала $j$, наблюдая сигнал
$j$ и рукопись, формирует апостериорное убеждение $\mu_{E_j}(\theta_A
\mid j)$ и принимает решение о рецензировании.

Ожидаемая полезность автора типа $\theta_A$ от подачи в журнал $j$:
\begin{equation}
  V_A(j \mid \theta_A) = p_j(\theta_A) \cdot q_j - c(j, \theta_A),
  \label{eq:vA}
\end{equation}
где $p_j(\theta_A)$~--- вероятность принятия, $c(j, \theta_A)$~---
ожидаемые издержки подачи (включая временные задержки, альтернативные
издержки отклонения и пересылки).

\subsection{Условие Спенса--Миррлиса и существование разделяющего PBE}

\begin{assumption}[Условие Спенса--Миррлиса]
  \label{assume:sc}
  Функция издержек $c(j, \theta_A)$ дважды непрерывно дифференцируема
  по обоим аргументам и удовлетворяет
  \begin{equation}
    \frac{\partial^2 c(j, \theta_A)}{\partial j \, \partial \theta_A}
    < 0 \quad \forall (j, \theta_A).
    \label{eq:single-crossing}
  \end{equation}
\end{assumption}

Условие~\eqref{eq:single-crossing} означает, что предельные издержки
подачи в более престижный журнал убывают с ростом качества автора:
высококачественным статьям относительно проще проходить жёсткое
рецензирование. Гладкое условие
Спенса--Миррлиса~\eqref{eq:single-crossing} достаточно, но не
эквивалентно порядковому свойству строгого пересечения
(single-crossing) в смысле Милгрома--Шеннона: первое сильнее и
исключает объединяющие (pooling) равновесия, тогда как второе более
общо и в принципе их допускает~\citep{edlin1998}. Для целей
построения разделяющего равновесия достаточно~\eqref{eq:single-crossing}.

\begin{theorem}[Существование разделяющего PBE]
  \label{thm:pbe-journal}
  В подыгре выбора журнала $\Gamma^{\mathrm{J}}$ при выполнении
  Предположения~\ref{assume:sc}, непрерывности $p_j(\theta_A)$ по
  $\theta_A$ и строгой монотонности $\kappa_j$ по $j$ существует
  разделяющее совершенное байесовское равновесие $(\sigma_A^{*},
  \mu^{*})$, в котором стратегия выбора журнала $j^{*}(\theta_A)$
  является монотонно возрастающей функцией типа автора, а
  апостериорные убеждения $\mu^{*}$ совместимы с байесовским
  обновлением на равновесной траектории.
\end{theorem}

\begin{proof}
  Определим семейство пороговых стратегий автора
  \[
    j^*(\theta_A) = k \iff \theta_A \in [\theta_{k-1}, \theta_k),
    \quad 0 = \theta_0 < \theta_1 < \dots < \theta_J = 1,
  \]
  где пороги $\{\theta_k\}$ подлежат определению. Редакторы при
  разделяющих сигналах формируют вырожденные апостериорные убеждения
  $\mu^{*}(\theta_A \mid j = k) = \mathbf{1}_{[\theta_{k-1}, \theta_k)}$.

  Условие безразличия автора типа $\theta_k$ между журналами $k$ и $k+1$:
  \begin{equation}
    V_A(k \mid \theta_k) = V_A(k+1 \mid \theta_k).
    \label{eq:indifference}
  \end{equation}
  По Предположению~\ref{assume:sc} функция $V_A(k+1 \mid \theta) -
  V_A(k \mid \theta)$ строго возрастает по $\theta$ (single-crossing
  эквивалентно монотонности разности). Следовательно,
  уравнение~\eqref{eq:indifference} имеет единственное решение
  $\theta_k^{*}$ на интервале $(\theta_{k-1}, 1)$. По индукции по $k
  = 1, \dots, J-1$ пороги $\theta_k^{*}$ однозначно определены.

  Проверим три требования PBE. (i)~Секвенциальная рациональность
  автора: по построению~\eqref{eq:indifference} автор типа $\theta
  \in [\theta_{k-1}^{*}, \theta_k^{*})$ получает максимум ожидаемой
  полезности при подаче в журнал $k$. (ii)~Секвенциальная
  рациональность редакторов: при вырожденных апостериорных убеждениях
  решение о рецензировании оптимально по определению.
  (iii)~Байесовское обновление: на равновесной траектории
  апостериорные убеждения согласованы с априорным распределением $F$
  по формуле Байеса.

  Для off-path сигналов применим критерий интуитивности
  Чо--Крепса~\citep{cho1987}: пусть $j'$~--- сигнал, не используемый
  ни одним типом в равновесии; назначим $\mu^{*}(\theta_A \mid j') =
  \delta_{\theta^*}$ для наибольшего типа $\theta^*$, для которого
  отклонение в $j'$ не доминируется равновесной стратегией. Такое
  назначение удовлетворяет интуитивному критерию и исключает часть
  неразумных равновесий.

  \begin{remark}
    При числе типов автора больше двух (а тем более при континууме
    типов $\theta_A \in [0,1]$) интуитивный критерий Чо--Крепса,
    вообще говоря, не выделяет единственное равновесие. Для уточнения
    off-path убеждений и отбора единственного разделяющего равновесия
    в этом случае применяется более сильный критерий
    D1~\citep{banks1987}, при котором off-path сигнал приписывается
    типу, имеющему наибольший относительный стимул к отклонению.
    Теорема~\ref{thm:pbe-journal} утверждает существование, но не
    единственность разделяющего PBE.
  \end{remark}

  Таким образом, $(\sigma_A^{*}, \mu^{*})$~--- разделяющее PBE.
  Отметим, что монотонные разделяющие равновесия в сигнальных играх
  восходят к~\citep{spence1973}; Теорема~\ref{thm:pbe-journal} есть
  адаптация этой классической конструкции к подыгре выбора журнала, а
  не самостоятельный результат теории сигналов.
\end{proof}

\begin{remark}
  В условиях российской системы журналов ранжирование $q_j$
  естественным образом индуцируется иерархией
  \[
    q_{\mathrm{WoS\,Q1}} > q_{\mathrm{RSCI\,K1}} >
    q_{\mathrm{\text{ядро\,РИНЦ}}} > q_{\mathrm{BAK\,K1}} >
    q_{\mathrm{K2}} > q_{\mathrm{K3}} > q_{\mathrm{\text{не-ВАК}}}.
  \]
  Пороги принятия $\kappa_j$ коррелируют с импакт-факторами и
  экспертными оценками~\citep{akoev2021, moskaleva2019}.
\end{remark}

% ============================================================
\section{Подыгра редактор-рецензент и принудительное цитирование}
\label{sec:coercive}

\subsection{Модель coercive citation}

Феномен принудительного цитирования, эмпирически
задокументированный~\citep{wilhite2012, heneberg2016}, формализуется
следующим образом. Редактор журнала $j$ при принятии рукописи имеет
право потребовать $k_j \in \N$ дополнительных ссылок на $j$. Такое
требование повышает двухлетний импакт-фактор $\mathrm{IF}_j$, но
несёт риск разоблачения с вероятностью $P_{\mathrm{det}} \in (0,1)$ и
репутационным штрафом $C_R > 0$ (исключение из баз, суспенсия
Clarivate, ретракция).

Ожидаемая полезность редактора от требования $k_j$ ссылок:
\begin{equation}
  U_{E_j}(k_j) = \beta_{\mathrm{IF}} \cdot \frac{\partial
  \mathrm{IF}_j}{\partial k_j} \cdot k_j - P_{\mathrm{det}}(k_j) \cdot C_R,
  \label{eq:editor-utility}
\end{equation}
где $\beta_{\mathrm{IF}} > 0$~--- институциональная ценность
импакт-фактора (финансирование, престиж), $P_{\mathrm{det}}(k_j)$~---
возрастающая функция числа требований.

\begin{proposition}[Равновесие подыгры принудительного цитирования]
  \label{prop:coercive}
  Рассмотрим двухходовую подыгру: редактор выбирает требование $k_j
  \geq 0$, затем автор выбирает между согласием (compliance) и
  отказом с пересылкой в другой журнал (reject). В совершенном по
  подыграм равновесии:
  \begin{enumerate}
    \item автор соглашается на $k_j$ принудительных ссылок тогда и
      только тогда, когда
      \begin{equation}
        V_A(\mathrm{accept}\mid k_j) - k_j \cdot r > V_A(\mathrm{reject}),
        \label{eq:author-accept}
      \end{equation}
      где $r > 0$~--- предельная репутационная издержка одного
      вынужденного цитирования;
    \item предвидя реакцию автора, редактор требует $k_j^* > 0$
      дополнительных ссылок тогда и только тогда, когда
      \begin{equation}
        \beta_{\mathrm{IF}} \cdot \frac{\partial
        \mathrm{IF}_j}{\partial k_j}\biggr|_{k_j=0} >
        P_{\mathrm{det}}'(0) \cdot C_R
        \label{eq:coercive-condition}
      \end{equation}
      и при этом ограничение участия автора~\eqref{eq:author-accept}
      не нарушено для выбранного $k_j^*$.
  \end{enumerate}
\end{proposition}

\begin{proof}
  Решаем подыгру обратной индукцией. На втором ходе автор сравнивает
  полезность согласия $V_A(\mathrm{accept}\mid k_j) - k_j r$ с
  полезностью отказа $V_A(\mathrm{reject})$,
  откуда~\eqref{eq:author-accept}. На первом ходе редактор
  максимизирует $U_{E_j}(k_j)$ из~\eqref{eq:editor-utility} при
  ограничении, что $k_j$ не превышает максимума, при котором автор
  ещё соглашается (из~\eqref{eq:author-accept}). Необходимое условие
  внутреннего максимума $U_{E_j}'(0) > 0$
  даёт~\eqref{eq:coercive-condition}. Структура
  условия~\eqref{eq:coercive-condition} есть стандартное беккеровское
  условие участия при правонарушении (предельная выгода против
  ожидаемого штрафа) и не претендует на новизну; содержательная
  часть~--- учёт реакции автора и ограничения участия, как в моделях
  редакционного процесса~\citep{garcia2015authored}.
\end{proof}

\subsection{Эмпирические следствия}

Предложение~\ref{prop:coercive} согласуется с эмпирическими
паттернами~\citep{wilhite2012}: около одной пятой опрошенных
академиков сталкивались с принудительным цитированием. Редакторы
стратегически таргетируют наиболее уязвимых авторов:
согласно~\citep{wilhite2012}, статус ассистент-профессора (в
сравнении с профессором) повышает вероятность подвергнуться
принуждению примерно на 5,5 процентного пункта, а каждый
дополнительный соавтор снижает её примерно на 2,0 процентного пункта.
В терминах модели это соответствует низкому $V_A(\mathrm{reject})$ у
младших авторов в~\eqref{eq:author-accept} и распределению
репутационной издержки $r$ между соавторами. Дальнейшее
теоретико-игровое моделирование редакционного процесса см.
в~\citep{garcia2015authored}.

% ============================================================
\section{Кооперативная подмодель соавторства и значение Шепли}
\label{sec:shapley}

\subsection{Характеристическая функция соавторской игры}

Для публикации $P$ с множеством потенциальных соавторов $N_P$
определим кооперативную игру с трансферабельной полезностью $(N_P,
v_P)$. Прежде чем задать $v_P$, выделим принципиальный
методологический момент. Карпов~\citep{karpov2014} доказал, что для
\emph{стандартной} задачи разделения соавторского кредита, в которой
публикация выступает неделимым ресурсом, а соавторы внутри одной
публикации взаимозаменяемы (анонимны в характеристической функции),
значение Шепли \emph{тождественно} равным долям $1/|N_P|$ независимо
от асимметрии продуктивности авторов по корпусу в целом.
Следовательно, чтобы метрика отклонялась от равных долей, необходимо
отказаться от внутрипубликационной анонимности и задать
характеристическую функцию, в которой предельные вклады соавторов
различаются \emph{внутри} самой публикации.

Введём поэтому вектор индивидуальных весов вкладов $\bm{w} = (w_i)_{i
\in N_P}$, $w_i > 0$, отражающих неравную роль соавторов в конкретной
работе $P$ (например, постановка задачи, основное доказательство,
эмпирическая часть). Характеристическая функция:
\begin{equation}
  v_P(S) = q_{j(P)} \cdot \Pr(\mathrm{accept} \mid S) \cdot \ind{S
  \supseteq K(P)} \cdot \Phi\!\left(\textstyle\sum_{i \in S} w_i\right),
  \label{eq:char-function}
\end{equation}
где $q_{j(P)}$~--- престиж целевого журнала, $K(P) \subseteq N_P$~---
критическое ядро соавторов (без которых публикация невозможна),
$\Phi\colon \R_+ \to \R_+$~--- возрастающая функция отдачи от
суммарного взвешенного вклада, эмпирически растущая с размером и
качеством команды~\citep{wuchty2007}. Зависимость $\Phi$ от
взвешенной суммы $\sum_{i \in S} w_i$, а не от мощности $|S|$,
нарушает внутрипубликационную анонимность и выводит игру из класса,
для которого справедлива эквивалентность Карпова. Такой подход
соответствует взвешенному значению Шепли~\citep{monderer1992}.

\subsection{Значение Шепли как распределение авторского кредита}

Значение Шепли соавтора $i \in N_P$:
\begin{equation}
  \varphi_i(v_P) = \sum_{S \subseteq N_P \setminus \{i\}} \frac{|S|!
  \cdot (|N_P| - |S| - 1)!}{|N_P|!} \cdot \bigl[v_P(S \cup \{i\}) -
  v_P(S)\bigr].
  \label{eq:shapley-def}
\end{equation}
Совокупный авторский кредит автора $A$ по публикациям $\mathcal{P}_A$:
\begin{equation}
  \mathrm{Credit}_A = \sum_{P \in \mathcal{P}_A} \varphi_A(v_P).
  \label{eq:credit}
\end{equation}

\begin{theorem}[Единственность значения Шепли и условие смещения равных долей]
  \label{thm:shapley}
  \textup{(i)} Значение Шепли~\eqref{eq:shapley-def} является
  единственным отображением $\varphi\colon \mathcal{G}_{N_P} \to
  \R^{N_P}$, удовлетворяющим аксиомам эффективности, симметрии,
  фиктивного игрока и аддитивности~\citep{shapley1953}. \textup{(ii)}
  \textup{(Эквивалентность Карпова.)} Если характеристическая функция
  не различает соавторов внутри публикации, то есть $v_P(S)$
  инвариантна относительно перестановок игроков и зависит лишь от
  того, какие публикации «покрыты» коалицией $S$, то значение Шепли
  тождественно равным долям: $\varphi_i(v_P) = v_P(N_P)/|N_P|$ для
  всех $i$, причём \emph{независимо} от асимметрии продуктивности
  авторов по корпусу~\citep{karpov2014}. \textup{(iii)} Если же веса
  вкладов $\bm{w}$ в~\eqref{eq:char-function} не равны между собой
  ($\exists\, i, k\colon w_i \neq w_k$), то при строго возрастающей
  $\Phi$ значение Шепли отклоняется от равных долей, и для соавтора
  $i \in K(P)$ с наибольшим весом выполнено $\varphi_i(v_P) >
  v_P(N_P)/|N_P|$, то есть критический соавтор недооценён в метрике
  равных долей.
\end{theorem}

\begin{proof}
  Пункт (i)~--- классический результат~\citep{shapley1953}. Пункт
  (ii) есть теорема об эквивалентности~\citep{karpov2014}: для игр, в
  которых соавторы одной публикации взаимозаменяемы (анонимны), все
  они являются симметричными игроками внутри каждого слагаемого
  характеристической функции, поэтому по аксиоме симметрии получают
  равные доли в каждой публикации; суммирование по $\mathcal{P}_A$
  даёт равенство $\varphi_i = v_P(N_P)/|N_P|$ без каких-либо
  предположений о глобальной симметрии. Пункт (iii): при неравных
  весах $\bm{w}$ функция $v_P$ зависит от состава коалиции через
  $\sum_{i \in S} w_i$, поэтому игроки с разными $w_i$ перестают быть
  симметричными. Для $i$ с наибольшим весом и любой коалиции $S
  \subseteq N_P \setminus \{i\}$, удовлетворяющей $S \cup \{i\}
  \supseteq K(P)$, маржинальный вклад
  \[
    v_P(S \cup \{i\}) - v_P(S) = q_{j(P)} \Pr(\mathrm{accept}\mid
    S\cup\{i\})\Bigl[\Phi\bigl(\textstyle\sum_{k\in S} w_k +
    w_i\bigr) - \Phi\bigl(\textstyle\sum_{k\in S} w_k\bigr)\Bigr]
  \]
  строго больше аналогичного вклада любого игрока с меньшим весом (в
  силу возрастания $\Phi$). Усреднение маржинальных вкладов по всем
  перестановкам в~\eqref{eq:shapley-def} даёт $\varphi_i(v_P) > v_P(N_P)/|N_P|$.
\end{proof}

\begin{remark}
  Пункты (ii)--(iii) проясняют распространённое заблуждение. Само по
  себе использование значения Шепли вместо равных долей \emph{не}
  меняет распределения кредита: при стандартной (анонимной)
  постановке они тождественны~\citep{karpov2014}. Содержательное
  отличие возникает только при явном моделировании неравных вкладов
  внутри публикации (веса $\bm{w}$), что и реализовано
  в~\eqref{eq:char-function}.
\end{remark}

\subsection{Ядро игры и устойчивость коалиции}

Ядро игры $(N_P, v_P)$:
\begin{equation}
  C(v_P) = \Bigl\{\bm{x} \in \R^{N_P} : \sum_{i \in N_P} x_i =
  v_P(N_P),\ \sum_{i \in S} x_i \geq v_P(S)\ \forall S \subseteq N_P \Bigr\}.
  \label{eq:core}
\end{equation}

\begin{theorem}[Непустота ядра и устойчивость коалиции]
  \label{thm:core}
  Пусть игра $(N_P, v_P)$ сбалансирована в смысле
  Бондаревой--Шепли~\citep{bondareva1963, shapley1967}, то есть для
  любого сбалансированного семейства коалиций $\mathcal{B}$ с весами
  $\{\lambda_S\}$ выполнено $\sum_{S \in \mathcal{B}} \lambda_S
  v_P(S) \leq v_P(N_P)$. Тогда $C(v_P) \neq \emptyset$, и коалиция
  $N_P$ устойчива в том смысле, что ни одна подкоалиция $S \subset
  N_P$ не может гарантированно улучшить выплаты своим членам
  отклонением от $v_P(N_P)$.
\end{theorem}

Доказательство есть классический результат теоремы Бондаревой--Шепли.

% ============================================================
\section{Повторяющаяся игра цитатных картелей и теорема Фолка}
\label{sec:citation-rings}

\subsection{Модель взаимного цитирования}

Цитатный картель (\emph{citation ring}) из $K$ игроков (авторов или
журналов) моделируется как бесконечно повторяющаяся игра с дисконтом
$\delta \in (0,1)$. На каждом шаге каждый из $K$ игроков решает:
цитировать ли остальных членов картеля (действие $C$) или нет
(действие $D$). Мгновенные выплаты: $v_{\mathrm{coop}}$ при взаимной
кооперации всех; $v_{\mathrm{dev}} > v_{\mathrm{coop}}$ при
одноразовом отклонении (игрок не цитирует, но получает цитирования);
$\underline{v} < v_{\mathrm{coop}}$~--- minmax-платёж.

\subsection{Trigger-стратегия и условие устойчивости}

Рассматривается grim-trigger стратегия: каждый игрок кооперирует до
первого отклонения любого из членов картеля; после отклонения все
переходят к минимаксному наказанию отклонившегося. Корректное
применение теоремы Фолка к игре $K$ игроков требует двух
дополнительных условий~\citep{fudenberg1986}: \textup{(а)} множество
допустимых индивидуально-рациональных выигрышей имеет полную
размерность (full dimensionality), либо выполнено более слабое
условие неэквивалентных полезностей (NEU); \textup{(б)} фаза
наказания совершенна по подыграм, то есть наказывающим игрокам
выгодно осуществлять наказание (для grim-trigger это обеспечивается
  тем, что после всеобщего перехода к $D$ повторение $D$ есть
статическое равновесие Нэша стадийной игры).

При выполнении условий (а)--(б) минимаксный платёж $\underline{v}$
есть нижняя граница индивидуально-рационального выигрыша. Условие,
что grim-trigger образует подыгровое совершенное равновесие для
каждого из $K$ игроков:
\begin{equation}
  \frac{v_{\mathrm{coop}}}{1 - \delta} \geq v_{\mathrm{dev}} +
  \frac{\delta}{1 - \delta} \underline{v},
  \label{eq:spe-ring}
\end{equation}
что эквивалентно
\begin{equation}
  \delta \geq \delta^{*} = \frac{v_{\mathrm{dev}} -
  v_{\mathrm{coop}}}{v_{\mathrm{dev}} - \underline{v}}.
  \label{eq:delta-star}
\end{equation}
Поскольку в симметричном картеле условие~\eqref{eq:spe-ring}
идентично для всех игроков, единый порог $\delta^*$ корректен; в
асимметричном случае порог определяется как максимум индивидуальных
порогов $\delta^*_i$.

\subsection{Верхняя граница размера устойчивого картеля}

Условие~\eqref{eq:delta-star} обеспечивает существование
кооперативного SPE, но \emph{само по себе} не ограничивает размер
картеля. Дополнительное ограничение возникает из растущей вероятности
разоблачения. Введём вероятность детекции картеля алгоритмами
выявления (например, CIDRE~\citep{kojaku2021}) как возрастающую
непрерывную функцию $P_{\mathrm{det}}(K)$ размера картеля $K$, и
ожидаемый штраф за разоблачение $C_{\mathrm{ring}} > 0$, вычитаемый
из дисконтированного кооперативного потока. Следующее утверждение
выводит границу размера картеля \emph{отдельно} от теоремы Фолка, как
самостоятельное следствие условия устойчивости при наличии детекции.

\begin{theorem}[Верхняя граница размера цитатного картеля при детекции]
  \label{thm:folk}
  Пусть выполнены условия \textup{(а)--(б)} теоремы Фолка, $\delta
  \geq \delta^*$, а $P_{\mathrm{det}}(K)$ строго возрастает и
  непрерывна. Тогда в повторяющейся игре цитатного картеля с
  детекцией устойчивый в SPE картель имеет размер $K \leq K^{*}$, где
  $K^{*}$~--- наибольшее $K$, удовлетворяющее неравенству
  \begin{equation}
    \frac{v_{\mathrm{coop}}(K) - P_{\mathrm{det}}(K) \cdot
    C_{\mathrm{ring}}}{1 - \delta} \geq v_{\mathrm{dev}} +
    \frac{\delta}{1 - \delta} \underline{v}.
    \label{eq:ring-bound}
  \end{equation}
\end{theorem}

\begin{proof}
  По теореме Фолка~\citep{fudenberg1986} при условиях (а)--(б) и
  $\delta \geq \delta^*$ кооперативный профиль поддерживается в SPE
  grim-trigger стратегией. Учёт детекции модифицирует ожидаемый
  кооперативный поток стадии: вместо $v_{\mathrm{coop}}(K)$ игрок
  получает $v_{\mathrm{coop}}(K) - P_{\mathrm{det}}(K) \cdot
  C_{\mathrm{ring}}$. Подставляя это в условие отсутствия выгоды от
  отклонения~\eqref{eq:spe-ring}, получаем~\eqref{eq:ring-bound}.
  Левая часть~\eqref{eq:ring-bound} строго убывает по $K$ (так как
    $P_{\mathrm{det}}$ строго возрастает, а $v_{\mathrm{coop}}(K)$
  растёт не быстрее линейного штрафа при больших $K$); поэтому
  существует единственное пороговое $K^*$, начиная с которого
  неравенство нарушается и кооперация перестаёт быть равновесной.
\end{proof}

\begin{corollary}
  Для типичных параметров российской экосистемы ($\delta \approx
    0{,}9$; $P_{\mathrm{det}}(K)$ оценённое по данным Clarivate о
    суспенсии 46 пар журналов и 55 изданий из JCR по причине citation
  stacking к 2019 году) $K^{*} \in \{2, 3, 4\}$. Это теоретически
  объясняет наблюдаемое доминирование малых циклов и распад больших картелей.
\end{corollary}

% ============================================================
\section{Эволюционная динамика стратегий в популяции}
\label{sec:ess}

\subsection{Популяционная модель}

Рассмотрим популяцию учёных с тремя чистыми стратегиями:
\begin{itemize}
  \item $s_1$ (глубинная): редкие высокоцитируемые публикации, низкое
    самоцитирование, выбор престижных журналов;
  \item $s_2$ (продуктивная): высокая публикационная активность,
    умеренное самоцитирование, диверсификация по журналам средней престижности;
  \item $s_3$ (картельная): систематическое самоцитирование, участие
    в цитатных картелях.
\end{itemize}

Пусть $\bm{x} = (x_1, x_2, x_3) \in \Delta^2$~--- вектор частот
стратегий в популяции. Матрица выплат $A = (a_{ij})_{i,j=1}^{3}$, где
$a_{ij}$~--- ожидаемый прирост h-индекса стратегии $i$ при встрече с
популяцией, играющей $j$.

\subsection{Репликативная динамика и ESS}

Репликативная динамика~\citep{taylor1978}:
\begin{equation}
  \dot{x}_i = x_i \bigl[(A\bm{x})_i - \bm{x}^{\top} A \bm{x}\bigr],
  \quad i = 1, 2, 3.
  \label{eq:replicator}
\end{equation}

\begin{definition}[ESS по Мэйнарду Смиту]
  Стратегия $\bm{x}^{*} \in \Delta^2$ является эволюционно устойчивой
  (ESS), если выполнены два условия:
  \begin{enumerate}
    \item $\bm{x}^{*\top} A \bm{x}^{*} \geq \bm{y}^{\top} A
      \bm{x}^{*}$ для всех $\bm{y} \in \Delta^2$;
    \item при равенстве в п.~1 выполнено $\bm{x}^{*\top} A \bm{y} >
      \bm{y}^{\top} A \bm{y}$ для всех $\bm{y} \neq \bm{x}^{*}$.
  \end{enumerate}
\end{definition}

\begin{remark}
  Существование ESS не является универсальным свойством матричных
  игр: игра $3 \times 3$ может вовсе не иметь ESS (классический
    пример~--- циклическая структура типа «камень-ножницы-бумага»,
    порождающая в~\eqref{eq:replicator} замкнутые орбиты вокруг
    внутренней неподвижной точки без асимптотической
  устойчивости)~\citep{hofbauer1998}. Поэтому формулируем
  существование ESS \emph{условно}, при явных ограничениях на матрицу
  выплат $A$.
\end{remark}

\begin{proposition}[Достаточное условие существования ESS]
  \label{prop:ess}
  Пусть существует внутренняя неподвижная точка $\bm{x}^{*} \in
  \mathrm{int}\,\Delta^2$ динамики~\eqref{eq:replicator}, и пусть
  матрица $A$ отрицательно определена на касательном подпространстве симплекса
  \[
    T\Delta^2 = \bigl\{\bm{\xi} \in \R^3 : \textstyle\sum_i \xi_i = 0\bigr\},
  \]
  то есть $\bm{\xi}^{\top} A \bm{\xi} < 0$ для всех $\bm{\xi} \in
  T\Delta^2 \setminus \{\bm{0}\}$. Тогда $\bm{x}^{*}$~---
  единственная ESS и глобально асимптотически устойчивая точка
  репликативной динамики~\eqref{eq:replicator} на $\mathrm{int}\,\Delta^2$.
\end{proposition}

\begin{proof}
  Условие отрицательной определённости $A$ на $T\Delta^2$
  эквивалентно тому, что $\bm{x}^{*}$ есть ESS~\citep{hofbauer1998}:
  для любой $\bm{y} \neq \bm{x}^{*}$, полагая $\bm{\xi} = \bm{y} -
  \bm{x}^{*} \in T\Delta^2$, получаем
  \[
    \bm{x}^{*\top} A \bm{y} - \bm{y}^{\top} A \bm{y} =
    -\,\bm{\xi}^{\top} A \bm{\xi} + (\bm{x}^{*} - \bm{y})^{\top} A \bm{x}^{*}.
  \]
  В силу того, что $\bm{x}^*$~--- неподвижная точка (выплаты всех
  присутствующих стратегий равны), второе слагаемое неотрицательно, а
  первое строго положительно по отрицательной определённости, откуда
  условие ESS выполнено. Глобальная асимптотическая устойчивость
  доказывается через функцию Ляпунова Кульбака--Лейблера $H(\bm{x}) =
  \sum_i x_i^{*} \ln(x_i^{*}/x_i)$, производная которой вдоль
  траекторий~\eqref{eq:replicator} равна $\dot{H} = -(\bm{x} -
  \bm{x}^{*})^{\top} A (\bm{x} - \bm{x}^{*}) > 0$ при $\bm{x} \neq
  \bm{x}^*$, что обеспечивает сходимость к $\bm{x}^{*}$~\citep{hofbauer1998}.
\end{proof}

\begin{remark}
  Эмпирически условие отрицательной определённости отвечает ситуации,
  когда ни одна чистая стратегия не доминирует: преимущество каждой
  стратегии убывает по мере роста её частоты (частотно-зависимый
  отбор с отрицательной обратной связью). Если же при малом
  $P_{\mathrm{det}}$ картельная стратегия $s_3$ строго доминирует
  (положительная обратная связь), внутренней ESS не существует, и
  динамика сходится к вершине $s_3$~--- вырожденному равновесию
  «поголовного картеля». Этот случай качественно соответствует
  феномену «естественного отбора плохой науки»~\citep{smaldino2016},
  который, однако, является агентной моделью культурной эволюции с
  отбором, а не доказательством ESS в смысле приведённого определения.
\end{remark}

% ============================================================
\section{Оптимальное самоцитирование: верхняя оценка прироста h-индекса}
\label{sec:selfcite}

\subsection{Постановка задачи}

Пусть текущий вектор цитирований автора $\bm{c} = (c_1, c_2, \dots,
c_n)$, упорядоченный по невозрастанию, $h = h(\bm{c})$. Автор имеет
бюджет $B \in \N$ дополнительных самоцитирований и распределяет их
вектором $\bm{s} = (s_1, \dots, s_n)$, $s_i \in \N_0$, $\sum_i s_i
\leq B$. Задача максимизации:
\begin{equation}
  \max_{\bm{s}} \Bigl[ h(\bm{c} + \bm{s}) - \lambda \cdot
  q_{\mathrm{susp}}(\bm{s}) \Bigr],
  \label{eq:selfcite-problem}
\end{equation}
где $q_{\mathrm{susp}}$~--- индекс подозрительности
(Bartneck--Kokkelmans q-index)~\citep{bartneck2011}, $\lambda > 0$.

\subsection{Пороговая структура h-индекса}

\begin{lemma}[Маргинальный прирост h]
  \label{lemma:marginal}
  Пусть $h = h(\bm{c})$, и статьи упорядочены по невозрастанию
  цитирований. Обозначим $m = |\{i : c_i \geq h+1\}|$~--- число
  статей, уже имеющих не менее $h+1$ цитирований; по определению
  h-индекса $m \leq h$. Тогда h-индекс возрастает до $h+1$ при
  добавлении самоцитирований тогда и только тогда, когда $(h+1)$-я по
  цитируемости статья достигает $h+1$ цитирований, что требует
  одновременного выполнения двух условий: \textup{(i)} имеется не
  менее $h+1-m$ статей с числом цитирований ровно $h$ (то есть в
  полуинтервале $[h, h+1)$), и \textup{(ii)} каждой из них добавлена
  хотя бы одна ссылка. Минимальный бюджет для прироста на единицу равен $h+1-m$.
\end{lemma}

\begin{proof}
  По определению~\eqref{eq:hindex-def} переход $h \to h+1$
  эквивалентен достижению $|\{i : c_i \geq h+1\}| \geq h+1$. Исходно
  это множество содержит $m \leq h$ статей. Добавление ссылок к
  статьям с $c_i > h$ не меняет их принадлежности множеству, к
  статьям с $c_i < h$~--- бесполезно (одна ссылка не доводит до
  $h+1$). Продуктивны только статьи с $c_i = h$: каждая такая статья
  при добавлении одной ссылки переходит в множество $\{c_i \geq
  h+1\}$. Чтобы мощность достигла $h+1$, нужно перевести туда $h+1-m$
  статей, что требует наличия не менее $h+1-m$ статей с $c_i = h$ и
  по одной ссылке на каждую. Отсюда минимальный бюджет $h+1-m$.
\end{proof}

\subsection{Основная оценка}

\begin{theorem}[Верхняя оценка прироста h от самоцитирования]
  \label{thm:selfcite}
  Для любого бюджета самоцитирований $B \geq 1$ и любой стратегии
  распределения $\bm{s}$ с $\sum_i s_i \leq B$ выполнена оценка
  \begin{equation}
    \Delta h(\bm{s}) = h(\bm{c} + \bm{s}) - h(\bm{c}) \leq
    \sum_{k=0}^{B-1} \ind{|\{i : c_i \in [h+k,\, h+k+1)\}| \geq 1}.
    \label{eq:selfcite-bound}
  \end{equation}
  В частности, прирост $\Delta h$ строго монотонно убывает по $B$ в
  смысле предельной отдачи: переход от прироста $k$ к $k+1$ требует
  наличия статей с $c_i = h+k$.
\end{theorem}

\begin{proof}
  Доказательство ведётся индукцией по $k = \Delta h$. База $k = 0$
  тривиальна. Переход: пусть прирост $\Delta h \geq k$ достигнут, и
  текущий h-индекс вспомогательного вектора $\bm{c} + \bm{s}$ равен
  $h + k$. По Лемме~\ref{lemma:marginal} прирост до $h + k + 1$
  возможен лишь при наличии хотя бы одной статьи с числом цитирований
  в полуинтервале $[h+k, h+k+1)$ и требует расходования части бюджета
  на доведение $(h+k+1)$-й статьи до порога. Если статей в этом
  полуинтервале нет, прирост невозможен независимо от оставшегося
  бюджета. Суммируя индикаторы наличия статей по уровням $k = 0,
  \dots, B-1$, получаем~\eqref{eq:selfcite-bound}.
\end{proof}

\begin{remark}
  Связь оценки~\eqref{eq:selfcite-bound} с убывающей отдачей
  самоцитирования установлена эмпирически и
  симуляционно~\citep{bartneck2011}: авторы способны заметно завысить
  h-индекс самоцитированиями, однако предельная отдача быстро падает,
  а сама стратегия выявляется q-индексом подозрительности. Наиболее
  результативная стратегия для устойчивого роста h~--- публикация
  работ, цитируемых другими, а не самоцитирование~\citep{bartneck2011}.
\end{remark}

\begin{corollary}[Закон убывающей полезности самоцитирования]
  \label{cor:diminishing}
  Эффективность бюджета $B$ на прирост h-индекса убывает: если
  множество «граничных» статей $B_h$ конечно, то существует $B^{*} =
  B^{*}(\bm{c})$ такое, что $\Delta h(\bm{s}) = \Delta h(\bm{s})|_{B
  = B^*}$ для всех $B > B^*$. Оптимальный бюджет определяется
  балансом прироста h-индекса и штрафа подозрительности
  в~\eqref{eq:selfcite-problem}.
\end{corollary}

\subsection{Эмпирический оптимум}

Численная калибровка Corollary~\ref{cor:diminishing} на типичных
распределениях цитирований российских авторов с учётом эмпирического
соотношения~\citep{fowler2007} (одно самоцитирование индуцирует
дополнительные внешние цитирования с коэффициентом $\sim 3$ за 5 лет)
даёт оптимальную долю самоцитирований в диапазоне 15--25\,\% от
общего числа цитирований автора. Это эмпирически проверяемое
предсказание модели.

% ============================================================
\section{Интеграция равновесий и импликации}
\label{sec:integration}

Совокупное стационарное равновесие
модели~\eqref{eq:game-structure}--\eqref{eq:utility-author}
объединяет разделяющее PBE Теоремы~\ref{thm:pbe-journal} в подыгре
выбора журнала, условное равновесие coercive citation
Предложения~\ref{prop:coercive}, распределение авторского кредита по
Шепли (Теорема~\ref{thm:shapley}), устойчивый размер цитатного
картеля Теоремы~\ref{thm:folk}, ESS популяционной стратегии
Предложения~\ref{prop:ess} и оптимальное самоцитирование
Теоремы~\ref{thm:selfcite}.

Модель даёт следующие проверяемые предсказания:
\begin{enumerate}
  \item[(P1)] Оптимальная доля самоцитирований автора лежит в
    интервале 15--25\,\% (Теорема~\ref{thm:selfcite},
    Следствие~\ref{cor:diminishing});
  \item[(P2)] Устойчивый размер цитатного картеля ограничен 2--4
    участниками (Теорема~\ref{thm:folk} и следствие);
  \item[(P3)] При неравных вкладах соавторов внутри публикации
    значение Шепли отклоняется от равных долей и недооценивает
    критических соавторов; в стандартной анонимной постановке
    отклонения нет (Теорема~\ref{thm:shapley}, пункты (ii)--(iii));
  \item[(P4)] Переход журнала между категориями ВАК (К2 $\to$ К1)
    индуцирует скачок среднего числа цитирований на статью (следствие
    разделяющего PBE и сигнальной структуры);
  \item[(P5)] Доля «продуктивной» стратегии $s_2$ в российской
    популяции превышает аналогичный показатель в международных
    выборках (Предложение~\ref{prop:ess});
  \item[(P6)] Авторы с однократно высоким h-индексом демонстрируют
    регрессию к среднему в последующие 5 лет (следствие impermanent
    reputation, приложимо к стохастическому потоку цитирований).
\end{enumerate}

Эмпирическая проверка предсказаний (P1)--(P6) составляет предмет
второй части исследования.

% ============================================================
\section{Заключение}
\label{sec:conclusion}

В работе построена единая многоагентная байесовская теоретико-игровая
модель научного публикования, в которой целевая функция задана как
взвешенная комбинация четырёх вариантов h-индекса РИНЦ. Модель
интегрирует сигнальную подыгру выбора журнала, подыгру coercive
citation, кооперативный аспект соавторства через значение Шепли,
повторяющуюся игру цитатных картелей и эволюционную динамику
стратегий в популяции.

Основные формальные результаты включают теорему существования
разделяющего PBE в подыгре выбора журнала
(Теорема~\ref{thm:pbe-journal}), условие равновесия coercive citation
(Предложение~\ref{prop:coercive}), теорему единственности значения
Шепли с уточнением условий отклонения от равных долей
(Теорема~\ref{thm:shapley}), верхнюю границу устойчивого размера
цитатного картеля (Теорема~\ref{thm:folk}), условия существования ESS
при параметрах российской системы (Предложение~\ref{prop:ess}) и
верхнюю оценку прироста h-индекса от стратегического самоцитирования
(Теорема~\ref{thm:selfcite}).

Модель приводит к шести проверяемым предсказаниям, эмпирическая
верификация которых на данных РИНЦ, Перечня ВАК и Белого списка РЦНИ
составляет содержание второй части исследования. Потенциальные
направления дальнейшего развития модели включают: динамическое
расширение с эндогенным входом новых журналов; многопериодную версию
с накоплением репутации и импермантной структурой убеждений
по~\citep{cripps2004}; расширение механизм-дизайна peer review в
духе~\citep{alon2011} с учётом импарциальности; сетевую модель
цитирования с эндогенным формированием связей.

\bibliographystyle{plainnat}
\bibliography{references}

\end{document}
